\section*{Chapitre \ref{ch:m2:inflation-markets} --- Marchés de l'inflation}

\begin{solution}{exo:m2:inflation-markets:1}
$292.296 + \tfrac{14}{31}\times(296.311 - 292.296) = 294.10922$~: août compte 31
jours, et le 15 correspond au jour 15, soit quatorze trente-et-unièmes du chemin parcouru.
\end{solution}

\begin{solution}{exo:m2:inflation-markets:2}
Exactement $1.0267/1.0014 - 1 = 2.5265\%$~; l'écart, 2.53~\%, le surestime
d'une fraction de point de base à ces niveaux.
\end{solution}

\begin{solution}{exo:m2:inflation-markets:3}
Le TIPS rembourse 100, son \omterm{def:m2:inflation-markets:linker}{plancher de déflation}~; le gilt rembourse 97, puisque son
principal suit l'indice à la baisse. Les coupons versés entre-temps portaient sur le
principal ajusté dans les deux cas.
\end{solution}

\begin{solution}{exo:m2:inflation-markets:4}
$100 \times (345/300 - 1.025^5) = \text{USD}~1.86$ million au fonds. Aucun à
l'indice à terme, $300 \times 1.025^5 = 339.42$.
\end{solution}

\begin{solution}{exo:m2:inflation-markets:5}
$294.10922/277.20274 = 1.06099$. Le 15 juillet, l'indice de référence interpolé
entre les indices d'avril et de mai, donc inférieur à ceux de mai et de juin alors que les prix
montaient~: le rapport valait 1.04814.
\end{solution}

\begin{solution}{exo:m2:inflation-markets:6}
$-0.144 - 0.359 - 0.349 = -0.85$ points de pourcentage sur les trois mois, soit environ
3.4 points par an en taux annualisé. Un \omterm{def:m2:inflation-markets:breakeven}{point mort d'inflation} lu sur une telle obligation sans
ajustement paraît 3.4 points trop bas.
\end{solution}

\begin{solution}{exo:m2:inflation-markets:7}
Sur 2010--2019, mars reste le mois le plus fort (0.355) et novembre le
plus faible ($-0.376$)~; aucun mois ne varie de plus de 0.11 point de pourcentage par rapport à
l'estimation 2010--2024, le plus grand écart concernant juin. Le profil est stable~; son
ampleur est estimée à environ un dixième de point.
\end{solution}

\begin{solution}{exo:m2:inflation-markets:8}
Le 9~\% représente l'inflation sur la dernière année~; le point mort est le prix de
l'inflation moyenne des dix prochaines années, avec les primes de risque et de liquidité incluses. La
position perçoit l'inflation réalisée contre le point mort seulement pendant les mois
où elle est détenue, et dans l'intervalle elle perd ou gagne selon le mouvement du point mort~; les
chiffres élevés déjà publiés sont dans le prix, à travers le couru connu.
La stratégie traduit l'anticipation d'une hausse du point mort, ou que l'inflation réalisée
sur la période de détention dépassera ce qui était anticipé, et non un simple encaissement.
\end{solution}

\begin{solution}{pb:m2:inflation-markets:1}
\textbf{1.} 292.296 (l'indice de mai) et 296.311 (juin).
\textbf{2.} 1.05445, le rapport de 292.296 à 277.20274.
\textbf{3.} 1.374~\%, la croissance de 292.296 à 296.311.
\textbf{4.} Le 13 juillet, lorsque l'indice de juin a été publié~: sur août, l'indice
de référence passe de l'indice de mai à celui de juin, tous deux alors connus.
\textbf{5.} 294.10922.
\textbf{6.} Un prix réel de 99.865 (sale) pour 100 de principal ajusté~; valeur de marché
$100 \times 0.99865 \times 1.05445$, USD~105.30 millions.
\textbf{7.} $105.30 \times 1.374\%$, USD~1\,446\,444.
\textbf{8.} USD~12\,513 (0.14~\% par an pendant 31 jours), plus un terme croisé de USD~172.
\textbf{9.} $105.30 \times 2.30\% \times 31/360$, USD~208\,558.
\textbf{10.} USD~1\,250\,572.
\textbf{11.} USD~27\,366~: le rendement de l'obligation nominale couvre à peine le repo.
\textbf{12.} 17.8~\% par an, contre un point mort de 2.53~\%~: le carry reflète la
publication d'un seul mois, et non l'anticipation du marché pour dix ans.
\textbf{13.} $296.276/296.311 - 1 = -0.012\%$~; carry $-202\,160$ dollars, le
coût du repo, sans rien pour le compenser.
\textbf{14.} USD~98\,921 par point de base~; 12.6 points de base de \omterm{def:m2:inflation-markets:breakeven}{rendement réel} suffiraient à
effacer le carry du mois.
\textbf{15.} Le couru est le même par dollar de valeur quelle que soit l'échéance, alors que
le risque de \omterm{def:m2:inflation-markets:breakeven}{rendement réel} croît avec la duration~: un TIPS à deux ans rapporte le même
couru d'août avec un cinquième du risque de prix.
\textbf{16.} Pas en soi~: le prix à terme pour la fin d'août inclut le
couru connu, si bien qu'un acheteur le 29 juillet le paie. Ce n'est un gain que face à
une position qui ne le détenait pas, ou si le \omterm{def:m2:inflation-markets:breakeven}{rendement réel} baisse.
\textbf{17.} Les \omterm{def:m2:inflation-markets:breakeven}{rendements réels} ont monté de deux points et demi, et une duration réelle au-delà
de neuf ans a transformé cela en une perte de prix plus forte que l'indexation de l'année.
\textbf{18.} Très peu~: le ratio d'indexation valait 1.054, si bien que l'indice devrait baisser
d'environ 5~\% d'ici 2032 avant que le plancher ne verse quoi que ce soit.
\textbf{19.} Résultat nommé~: \emph{le carry de l'\omterm{def:m2:inflation-markets:linker}{obligation indexée}} d'un nominal de USD~100 millions du
TIPS à dix ans sur août 2022 s'élevait à environ USD~1.25 million, connu le 13 juillet.
\textbf{20.} Une \omterm{def:m2:inflation-markets:linker}{obligation indexée} protège son détenteur contre le niveau des prix, et non contre
le taux d'intérêt réel.
\end{solution}

\begin{solution}{iq:m2:inflation-markets:1}
Le taux d'inflation constant auquel une obligation nominale et une \omterm{def:m2:inflation-markets:linker}{obligation indexée}
de même échéance rapportent le même taux, à peu près le rendement nominal moins le
\omterm{def:m2:inflation-markets:breakeven}{rendement réel}. C'est un prix~: il intègre une prime pour porter le risque d'inflation,
en retranchant la prime que doit verser l'\omterm{def:m2:inflation-markets:linker}{obligation indexée}, moins liquide, plus les effets de convexité, et il
varie avec les \omterm{def:m2:the-treasury-market:dash}{ruées vers la liquidité}, comme en mars 2020.
\iqlookfor{la relation de Fisher, et les primes qui la séparent d'une prévision.}
\end{solution}

\begin{solution}{iq:m2:inflation-markets:2}
Le principal est multiplié par le ratio d'indexation, c'est-à-dire l'indice de référence (le
CPI-U non désaisonnalisé, avec un décalage de trois mois, interpolé quotidiennement) rapporté
à sa valeur à la date considérée~; les coupons sont le taux réel fixe appliqué à ce principal.
À l'échéance, le détenteur reçoit le plus élevé entre le principal ajusté et le principal
initial~: la déflation sur toute la durée de vie est bornée par un plancher, mais les coupons
versés en cours de route sur un principal plus faible ne le sont pas.
\iqlookfor{l'indice décalé et interpolé, et le plancher portant uniquement sur le principal.}
\end{solution}

\begin{solution}{iq:m2:inflation-markets:3}
Parce que le principal croît sur un mois de la variation de l'indice de deux et trois
mois plus tôt, si bien qu'une forte publication se transforme en un couru important et connu quelques
semaines plus tard~; financée au \omterm{def:m2:repo-and-specials:rate}{taux de repo}, la position porte le couru moins le
repo. Le carry revient à celui qui détenait l'obligation avant que la publication soit intégrée au prix~; un acheteur
après la publication le paie dans le prix à terme.
\iqlookfor{le mécanisme du décalage, et le fait que le carry connu est intégré au prix.}
\end{solution}

\begin{solution}{iq:m2:inflation-markets:4}
Instruments et coûts différents~: le swap ne demande aucun financement et consomme peu de
bilan~; l'obligation doit être financée et mobilise du bilan~; leur liquidité
diffère~; le \omterm{def:m2:inflation-markets:linker}{plancher de déflation} de l'obligation a une valeur que le swap n'a pas~; l'offre et la demande
diffèrent (fonds de pension receveurs dans les swaps, émissions du Trésor en obligations)~; les
conventions de décalage et les dates exactes des indices peuvent ne pas coïncider~; et le
\omterm{def:m2:inflation-markets:breakeven}{point mort d'inflation} de l'obligation est coté en rendement, celui du swap en \omterm{def:m2:government-bonds:zero}{taux zéro-coupon}.
\iqlookfor{financement et bilan, le plancher, les flux et les conventions.}
\end{solution}

\begin{solution}{iq:m2:inflation-markets:5}
Estimer les coefficients mensuels à partir de l'historique de l'indice non corrigé (par exemple,
la variation de chaque mois moins la moyenne de son année, moyennée sur les années complètes, ou les
coefficients saisonniers officiels), vérifier leur stabilité sur des sous-périodes, les normaliser
pour que leur somme soit nulle sur l'année, et les appliquer à la courbe d'indice à terme
construite à partir des swaps, afin que les variations mensuelles à terme portent la structure du calendrier~;
puis évaluer les \omterm{def:m2:inflation-markets:linker}{linkers} courts sur cette courbe.
\iqlookfor{estimation, stabilité, normalisation et application à la courbe à terme.}
\end{solution}

\begin{solution}{iq:m2:inflation-markets:6}
Un tableau de conventions par obligation (indice, décalage, interpolation, arrondi, date de référence,
plancher, dates de coupon), les gilts à décalage de huit mois faisant l'objet d'une règle distincte~; un magasin
d'indices tenu par mois, alimenté par chaque institut statistique le jour de sa publication,
avec des contrôles par rapport à la valeur précédente, à son calendrier de publication et au
versionnement des rebasages~; une fonction pure allant de (obligation, date, magasin d'indices) vers le ratio,
testée sur les exemples publiés par les émetteurs et sur les tables de ratios~; une sortie qui est
réconciliée chaque matin avec les ratios publiés par les émetteurs eux-mêmes, et des alertes
se déclenchent lorsqu'un indice est en retard ou qu'une règle d'obligation manque.
\iqlookfor{conventions traitées comme des données, magasin d'indices versionné, et réconciliation avec
les tables des émetteurs.}
\end{solution}
